% 复平面：z = 3 + 4i，标 z, \bar z, |z|=5
\documentclass[tikz,border=4pt]{standalone}
\usepackage{ctex}
\usepackage{amsmath}
\usetikzlibrary{calc, arrows.meta}
\begin{document}
\begin{tikzpicture}[scale=0.8, thick]
  % 坐标轴
  \draw[->, gray!60!black] (-1, 0) -- (5.5, 0) node[right, black] {\small 实轴 $\mathrm{Re}$};
  \draw[->, gray!60!black] (0, -5.5) -- (0, 5.5) node[above, black] {\small 虚轴 $\mathrm{Im}$};

  % 网格
  \draw[gray!25, very thin] (-0.5,-5) grid (5, 5);
  \foreach \x in {1,2,3,4,5} \draw[gray] (\x, 0.08) -- (\x, -0.08) node[below, font=\scriptsize, black] {$\x$};
  \foreach \y/\label in {1/1,2/2,3/3,4/4,-1/{-1},-2/{-2},-3/{-3},-4/{-4}}
    \draw[gray] (0.08, \y) -- (-0.08, \y) node[left, font=\scriptsize, black] {$\label\,i$};

  % 原点
  \node[below left, font=\scriptsize] at (0,0) {$O$};

  % z = 3 + 4i 与 \bar z = 3 - 4i
  \coordinate (Z) at (3, 4);
  \coordinate (Zb) at (3, -4);

  % 向量
  \draw[->, blue!70!black, very thick] (0,0) -- (Z);
  \draw[->, red!70!black, very thick] (0,0) -- (Zb);

  % 点
  \fill[blue!70!black] (Z) circle (2pt);
  \fill[red!70!black] (Zb) circle (2pt);

  % 标注 z 和 \bar z
  \node[above right, blue!70!black] at (Z) {$z = 3+4i$};
  \node[below right, red!70!black] at (Zb) {$\overline{z} = 3-4i$};

  % |z| = 5 标注（向量长度）
  \node[blue!70!black, font=\small] at (1.0, 2.4) {$|z|=5$};

  % 辅助：实部、虚部分量
  \draw[gray, dashed] (Z) -- (3, 0);
  \draw[gray, dashed] (Z) -- (0, 4);
  \node[gray!50!black, font=\scriptsize, below] at (1.5, 0) {实部 $a=3$};
  \node[gray!50!black, font=\scriptsize, left] at (0, 2) {虚部 $b=4$};

  % 实轴对称提示
  \draw[gray, dashed] (Z) -- (Zb);
  \node[gray!50!black, font=\scriptsize, align=center] at (4.6, 0.4) {关于\\实轴对称};
\end{tikzpicture}
\end{document}
